\documentclass[10pt,a4paper]{amsart}
\usepackage{amsmath,amssymb,amsthm,graphicx}
% Retained minimal preamble for the Stacks Project chapter flat.tex (More on Flatness).
% Verbatim excerpt of the chapter preamble: Xy-pic setup (native preamble lines 14-18) and the native
% theorem declarations with the mathematical macros (native preamble lines 154-217). Copied mechanically;
% nothing added, removed or redefined. The wrapper class and the amsmath/amssymb/amsthm/graphicx base
% packages are supplied by the assembler, so the class/externaldocument part of the native preamble is omitted.
% For commutative diagrams we use Xy-pic
\usepackage[all]{xy}

% We use 2cell for 2-commutative diagrams.
\xyoption{2cell}
% Theorem environments.
%
\theoremstyle{plain}
\newtheorem{theorem}[subsection]{Theorem}
\newtheorem{proposition}[subsection]{Proposition}
\newtheorem{lemma}[subsection]{Lemma}

\theoremstyle{definition}
\newtheorem{definition}[subsection]{Definition}
\newtheorem{example}[subsection]{Example}
\newtheorem{exercise}[subsection]{Exercise}
\newtheorem{situation}[subsection]{Situation}

\theoremstyle{remark}
\newtheorem{remark}[subsection]{Remark}
\newtheorem{remarks}[subsection]{Remarks}

\numberwithin{equation}{subsection}

% Macros
%
\def\lim{\mathop{\mathrm{lim}}\nolimits}
\def\colim{\mathop{\mathrm{colim}}\nolimits}
\def\Spec{\mathop{\mathrm{Spec}}}
\def\Hom{\mathop{\mathrm{Hom}}\nolimits}
\def\Ext{\mathop{\mathrm{Ext}}\nolimits}
\def\SheafHom{\mathop{\mathcal{H}\!\mathit{om}}\nolimits}
\def\SheafExt{\mathop{\mathcal{E}\!\mathit{xt}}\nolimits}
\def\Sch{\mathit{Sch}}
\def\Mor{\mathop{\mathrm{Mor}}\nolimits}
\def\Ob{\mathop{\mathrm{Ob}}\nolimits}
\def\Sh{\mathop{\mathit{Sh}}\nolimits}
\def\NL{\mathop{N\!L}\nolimits}
\def\CH{\mathop{\mathrm{CH}}\nolimits}
\def\proetale{{pro\text{-}\acute{e}tale}}
\def\etale{{\acute{e}tale}}
\def\QCoh{\mathit{QCoh}}
\def\Ker{\mathop{\mathrm{Ker}}}
\def\Im{\mathop{\mathrm{Im}}}
\def\Coker{\mathop{\mathrm{Coker}}}
\def\Coim{\mathop{\mathrm{Coim}}}

% Boxtimes
%
\DeclareMathSymbol{\boxtimes}{\mathbin}{AMSa}{"02}

%
% Macros for moduli stacks/spaces
%
\def\QCohstack{\mathcal{QC}\!\mathit{oh}}
\def\Cohstack{\mathcal{C}\!\mathit{oh}}
\def\Spacesstack{\mathcal{S}\!\mathit{paces}}
\def\Quotfunctor{\mathrm{Quot}}
\def\Hilbfunctor{\mathrm{Hilb}}
\def\Curvesstack{\mathcal{C}\!\mathit{urves}}
\def\Polarizedstack{\mathcal{P}\!\mathit{olarized}}
\def\Complexesstack{\mathcal{C}\!\mathit{omplexes}}
% \Pic is the operator that assigns to X its picard group, usage \Pic(X)
% \Picardstack_{X/B} denotes the Picard stack of X over B
% \Picardfunctor_{X/B} denotes the Picard functor of X over B
\def\Pic{\mathop{\mathrm{Pic}}\nolimits}
\def\Picardstack{\mathcal{P}\!\mathit{ic}}
\def\Picardfunctor{\mathrm{Pic}}
\def\Deformationcategory{\mathcal{D}\!\mathit{ef}}

% --- appended 2026-10-12 (Chief): the project own definitions of the reference/slogan/history
% display environments, copied mechanically from the shared Stacks preamble lines 8-12
% (\usepackage{verbatim} plus the three \newenvironment lines, which the chapter preamble keeps
% verbatim from the project preamble). Nothing added, removed or redefined: these environments
% carry no mathematical content and the project itself typesets them as comments. Without these
% three lines a counted statement whose own text uses them fails with Environment slogan undefined.
\usepackage{verbatim}
\newenvironment{reference}{\comment}{\endcomment}
\newenvironment{slogan}{\comment}{\endcomment}
\newenvironment{history}{\comment}{\endcomment}
\usepackage{cleveref}
\crefname{originallemma}{original lemma}{original lemmas}
\Crefname{originallemma}{Original lemma}{Original lemmas}
\makeatletter
\newlabel{algebra-lemma-local-flat-ff}{{00HR}{source}{}{sourcepointer-1}{}}
\newlabel{algebra-lemma-local-flat-ff@cref}{{[originallemma][0][]00HR}{[1][0][]source}{}{}{}}
\makeatother
\crefname{originallemma}{original lemma}{original lemmas}
\Crefname{originallemma}{Original lemma}{Original lemmas}
\makeatletter
\newlabel{algebra-lemma-module-map-property-in-colimit}{{05LI}{source}{}{sourcepointer-2}{}}
\newlabel{algebra-lemma-module-map-property-in-colimit@cref}{{[originallemma][0][]05LI}{[1][0][]source}{}{}{}}
\makeatother
\crefname{originallemma}{original lemma}{original lemmas}
\Crefname{originallemma}{Original lemma}{Original lemmas}
\makeatletter
\newlabel{properties-lemma-finite-presentation-module}{{01PC}{source}{}{sourcepointer-3}{}}
\newlabel{properties-lemma-finite-presentation-module@cref}{{[originallemma][0][]01PC}{[1][0][]source}{}{}{}}
\makeatother
\crefname{originallemma}{original lemma}{original lemmas}
\Crefname{originallemma}{Original lemma}{Original lemmas}
\makeatletter
\newlabel{properties-lemma-finite-type-module}{{01PB}{source}{}{sourcepointer-4}{}}
\newlabel{properties-lemma-finite-type-module@cref}{{[originallemma][0][]01PB}{[1][0][]source}{}{}{}}
\makeatother
\crefname{originallemma}{original lemma}{original lemmas}
\Crefname{originallemma}{Original lemma}{Original lemmas}
\makeatletter
\newlabel{situation-iso}{{05MH}{source}{}{sourcepointer-5}{}}
\newlabel{situation-iso@cref}{{[originallemma][0][]05MH}{[1][0][]source}{}{}{}}
\makeatother
\crefname{originallemma}{original lemma}{original lemmas}
\Crefname{originallemma}{Original lemma}{Original lemmas}
\makeatletter
\newlabel{topologies-lemma-fpqc}{{022D}{source}{}{sourcepointer-6}{}}
\newlabel{topologies-lemma-fpqc@cref}{{[originallemma][0][]022D}{[1][0][]source}{}{}{}}
\makeatother
\title{Sheaf and limit preserving properties of the functors of points where a module map is injective, surjective, bijective or zero}
\author{The Stacks Project Authors}
\date{stacks-project:flat.tex:local-snapshot-20261011; GNU FDL 1.2 or later}
\begin{document}
\maketitle
\noindent\textit{Editorial note.} All delivered mathematical text is the author's own, copied from the pinned Stacks Project chapter file flat.tex; no statement or proof step was rewritten, reordered or removed.\par\medskip
\noindent\textit{Editorial note.} Preamble provenance: the chapter's own retained minimal preamble (sources/stacks-flat/retained-preamble.tex, lines 1-85) is a verbatim excerpt of the native Xy-pic setup and of the native theorem declarations and mathematical macros, followed by the project's own reference/slogan/history display-environment definitions copied mechanically from the shared preamble; the wrapper class and the amsmath/amssymb/amsthm/graphicx base packages are supplied by the assembler, and extra\_packages is left empty because everything needed is already in the retained preamble. No macro was invented or redefined.\par\medskip
\noindent\textit{Editorial note.} Same-chapter pointers are located at their \textbackslash{}label lines inside the pinned chapter file; cross-chapter pointers are located in the project's own tag list staged in the snapshot (sources/stacks-chow/tags.txt), at the verbatim lines carrying their canonical tags.\par\medskip
\noindent\textit{Allowed standard prerequisites (editorial).}\begin{itemize}
\item The situation and functors recording injectivity, surjectivity, bijectivity and vanishing of a module map. Same chapter flat.tex, situation labelled situation-iso (statement, proof line 5920). Fixes the four functors F\_iso, F\_inj, F\_surj and F\_zero whose sheaf and limit preserving properties are asserted.
\item Flatness of an fpqc covering and faithful flatness of the induced local maps. Topologies, label topologies-lemma-fpqc (proof line 5937); Algebra, label algebra-lemma-local-flat-ff (proof line 5942). Shows that the pullback covering of X\_T is an fpqc covering, so every point has a point of some member over it, and that the resulting local map of stalks is faithfully flat, which is what makes injectivity, surjectivity, bijectivity and vanishing equivalent upstairs and downstairs.
\item The algebraic colimit statements for module maps and reduction of finite type and finite presentation modules. Algebra, label algebra-lemma-module-map-property-in-colimit, parts (1), (2) and (3), invoked at proof lines 5975, 5987 and 6004; Properties, labels properties-lemma-finite-type-module (proof line 5966) and properties-lemma-finite-presentation-module (proof line 5994). Each of the limit preserving parts is reduced to the corresponding algebraic statement that surjectivity, vanishing or bijectivity of a base changed map descends to some stage of a directed colimit, and the reduction from sheaves to modules uses the affine local description of a finite type or finite presentation module.
\end{itemize}\medskip
\ifdefined\hypertarget\hypertarget{sourcepointer-1}{}\fi
\noindent\textit{Source reference.} Original lemma 00HR: sources/stacks-chow/tags.txt, source lines 653--653. Navigation pointer for \textbackslash{}ref\{algebra-lemma-local-flat-ff\} used inside the counted span (line 5942); the argument that consumes it is retained in the span.\par\medskip
\ifdefined\hypertarget\hypertarget{sourcepointer-2}{}\fi
\noindent\textit{Source reference.} Original lemma 05LI: sources/stacks-chow/tags.txt, source lines 7237--7237. Navigation pointer for \textbackslash{}ref\{algebra-lemma-module-map-property-in-colimit\} used inside the counted span (line 5975); the argument that consumes it is retained in the span.\par\medskip
\ifdefined\hypertarget\hypertarget{sourcepointer-3}{}\fi
\noindent\textit{Source reference.} Original lemma 01PC: sources/stacks-chow/tags.txt, source lines 2225--2225. Navigation pointer for \textbackslash{}ref\{properties-lemma-finite-presentation-module\} used inside the counted span (line 5994); the argument that consumes it is retained in the span.\par\medskip
\ifdefined\hypertarget\hypertarget{sourcepointer-4}{}\fi
\noindent\textit{Source reference.} Original lemma 01PB: sources/stacks-chow/tags.txt, source lines 2224--2224. Navigation pointer for \textbackslash{}ref\{properties-lemma-finite-type-module\} used inside the counted span (line 5966); the argument that consumes it is retained in the span.\par\medskip
\ifdefined\hypertarget\hypertarget{sourcepointer-5}{}\fi
\noindent\textit{Source reference.} Original lemma 05MH: sources/stacks-flat/pinned-original-source.tex, source lines 5896--5896. Navigation pointer for \textbackslash{}ref\{situation-iso\} used inside the counted span (line 5920); the argument that consumes it is retained in the span.\par\medskip
\ifdefined\hypertarget\hypertarget{sourcepointer-6}{}\fi
\noindent\textit{Source reference.} Original lemma 022D: sources/stacks-chow/tags.txt, source lines 2702--2702. Navigation pointer for \textbackslash{}ref\{topologies-lemma-fpqc\} used inside the counted span (line 5937); the argument that consumes it is retained in the span.\par\medskip
\begin{lemma}
\label{lemma-iso-sheaf}
In Situation \ref{situation-iso}.
\begin{enumerate}
\item Each of the functors $F_{iso}$, $F_{inj}$, $F_{surj}$, $F_{zero}$
satisfies the sheaf property for the fpqc topology.
\item If $f$ is quasi-compact and $\mathcal{G}$ is of finite type,
then $F_{surj}$ is limit preserving.
\item If $f$ is quasi-compact and $\mathcal{F}$ of finite type, then
$F_{zero}$ is limit preserving.
\item If $f$ is quasi-compact, $\mathcal{F}$ is of finite type, and
$\mathcal{G}$ is of finite presentation, then $F_{iso}$ is limit preserving.
\end{enumerate}
\end{lemma}

\begin{proof}
Let $\{T_i \to T\}_{i \in I}$ be an fpqc covering of schemes over $S$.
Set $X_i = X_{T_i} = X \times_S T_i$ and $u_i = u_{T_i}$.
Note that $\{X_i \to X_T\}_{i \in I}$ is an fpqc covering of $X_T$, see
Topologies, Lemma \ref{topologies-lemma-fpqc}.
In particular, for every $x \in X_T$ there exists an $i \in I$
and an $x_i \in X_i$ mapping to $x$. Since
$\mathcal{O}_{X_T, x} \to \mathcal{O}_{X_i, x_i}$ is flat, hence
faithfully flat (see
Algebra, Lemma \ref{algebra-lemma-local-flat-ff})
we conclude that $(u_i)_{x_i}$ is injective, surjective, bijective, or zero
if and only if $(u_T)_x$ is injective, surjective, bijective, or zero.
Whence part (1) of the lemma.

\medskip\noindent
Proof of (2). Assume $f$ quasi-compact and $\mathcal{G}$ of finite type.
Let $T = \lim_{i \in I} T_i$ be a directed limit of affine $S$-schemes
and assume that $u_T$ is surjective.
Set $X_i = X_{T_i} = X \times_S T_i$ and
$u_i = u_{T_i} : \mathcal{F}_i = \mathcal{F}_{T_i}
\to \mathcal{G}_i = \mathcal{G}_{T_i}$.
To prove part (2) we have to show that $u_i$ is surjective for some $i$.
Pick $i_0 \in I$ and replace $I$ by $\{i \mid i \geq i_0\}$.
Since $f$ is quasi-compact the scheme $X_{i_0}$ is quasi-compact.
Hence we may choose affine opens $W_1, \ldots, W_m \subset X$
and an affine open covering
$X_{i_0} = U_{1, i_0} \cup \ldots \cup U_{m, i_0}$ such that
$U_{j, i_0}$ maps into $W_j$ under the projection morphism $X_{i_0} \to X$.
For any $i \in I$ let $U_{j, i}$ be the inverse image of $U_{j, i_0}$.
Setting $U_j = \lim_i U_{j, i}$ we see that $X_T = U_1 \cup \ldots \cup U_m$
is an affine open covering of $X_T$. Now it suffices to show, for a given
$j \in \{1, \ldots, m\}$ that $u_i|_{U_{j, i}}$ is surjective for some
$i = i(j) \in I$. Using
Properties, Lemma \ref{properties-lemma-finite-type-module}
this translates into the following algebra problem:
Let $A$ be a ring and let $u : M \to N$ be an $A$-module map.
Suppose that $R = \colim_{i \in I} R_i$ is a directed colimit
of $A$-algebras. If $N$ is a finite $A$-module and if
$u \otimes 1 : M \otimes_A R \to N \otimes_A R$ is surjective, then
for some $i$ the map
$u \otimes 1 : M \otimes_A R_i \to N \otimes_A R_i$ is surjective.
This is
Algebra, Lemma \ref{algebra-lemma-module-map-property-in-colimit} part (2).

\medskip\noindent
Proof of (3). Exactly the same arguments as given in the proof of (2)
reduces this to the following algebra problem:
Let $A$ be a ring and let $u : M \to N$ be an $A$-module map.
Suppose that $R = \colim_{i \in I} R_i$ is a directed colimit
of $A$-algebras. If $M$ is a finite $A$-module and if
$u \otimes 1 : M \otimes_A R \to N \otimes_A R$ is zero, then
for some $i$ the map
$u \otimes 1 : M \otimes_A R_i \to N \otimes_A R_i$ is zero.
This is
Algebra, Lemma \ref{algebra-lemma-module-map-property-in-colimit} part (1).

\medskip\noindent
Proof of (4).
Assume $f$ quasi-compact and $\mathcal{F}, \mathcal{G}$ of finite presentation.
Arguing in exactly the same manner as in the previous paragraph
(using in addition also
Properties, Lemma \ref{properties-lemma-finite-presentation-module})
part (3) translates into the following algebra statement:
Let $A$ be a ring and let $u : M \to N$ be an $A$-module map.
Suppose that $R = \colim_{i \in I} R_i$ is a directed colimit
of $A$-algebras. Assume $M$ is a finite $A$-module, $N$ is a finitely
presented $A$-module, and
$u \otimes 1 : M \otimes_A R \to N \otimes_A R$ is an isomorphism.
Then for some $i$ the map
$u \otimes 1 : M \otimes_A R_i \to N \otimes_A R_i$ is an isomorphism.
This is
Algebra, Lemma \ref{algebra-lemma-module-map-property-in-colimit} part (3).
\end{proof}

\end{document}
